\section{A universe, mostly dark}
\label{sec:grdark}

Start from what the book can already do.  Given one expiry's quotes,
Part~I fits an arbitrage-free smile.  Part~II prepares the quotes that
feed it.  Chapters~9 and~10 gave each smile a life in time: a transport
rule for the day's spot move, and a filter that weighs each morning's
noisy reading against yesterday's transported estimate.  All of it is
\emph{per node}: one underlier, one expiry, its own quotes, its own
yesterday.

Now widen the lens.  A desk quotes options on hundreds of underliers,
each with a dozen expiries.  Call each (underlier, expiry) pair a
\emph{node}: one smile, one probability law.  On a normal morning a few
nodes are \emph{lit} --- the index expiries and a handful of large
names trade densely enough for a full, trustworthy calibration --- and
the rest are \emph{dark}: a thin name that has not traded today, a far
tenor with four stale quotes, a weekly that printed twice before
lunch.  The dark nodes still need surfaces.  Client quotes, risk
reports and scenario engines do not pause because the tape is thin.

The chapter's universe is drawn in \cref{fig:gruniverse}(a).  It is
built from the book's frozen data plus two synthetic names, and it is
small enough to read whole: twenty nodes in four rows.  The top two
rows are real --- the frozen SPY and NVDA boards, six expiries each,
from 18~days to 501~days.  (The two sub-week expiries of the frozen
boards are left out.  The next section will relate expiries through
the ratio of their maturities; across a two-day gap that ratio is
really a statement about scheduled events, and Chapter~8 taught which
clock such nodes actually live on.)  The bottom two rows are synthetic names built from
the real boards by stated recipes --- a \emph{sister} name constructed
beside NVDA, and a \emph{blend} constructed between the two boards ---
four middle expiries each; Appendix~\ref{app:grprotocol} gives the
recipes.  Filled orange nodes are lit this morning: the whole index
board, and two expiries of NVDA.  Open nodes are dark: \MacGrDarkCount\
of the \MacGrNodeCount, including the sister and blend boards in their
entirety.  The grey segments and arrows are the \emph{edges} --- who is
believed to inform whom --- and the rest of the chapter is about what,
exactly, an edge should say.

Panel~(b) shows the morning itself.  Each lit node was calibrated, and
each calibration is compared with what \cref{ch:infer} would have
predicted for that node --- yesterday's estimate, transported to today.
The gap is the \emph{innovation}, the same word and the same object as
in the filter update \cref{eq:infupdate}: the part of today that
yesterday did not predict.  This particular morning carries a broad
repricing, largest by far at the short end: the index printed
\MacGrSpyDecObs\ volatility points at the December node (137~days)
and fully \MacGrSpyShortObs\ at 18~days, with only its 410-day node
ticking slightly down; and NVDA's lit 46-day node printed
\MacGrNvdaSepObs\ --- a high-beta name amplifying the index's move at
that expiry.  Short-dated volatility swinging hardest in a general
repricing is exactly what Chapter~8's clock leads one to expect, and
\cref{sec:gredge} will turn that expectation into an amplitude.

\begin{figure}[htbp]
\centering
\paperfig{\textwidth}{fig_gr_universe.pdf}
\caption{The staged universe.  (a)~Twenty nodes in four rows: the real
SPY and NVDA boards (six expiries each) and two synthetic names built
from them.  Filled orange nodes are lit this morning
(\MacGrLitCount\ of \MacGrNodeCount); open nodes are dark.  Grey
segments are calendar edges within a name; arrows are cross-name
edges, pointing from the informing name to the informed one.
(b)~The eight lit innovations: each lit node's calibrated handle minus
its transported baseline, in volatility points.  The morning is a
broad repricing, largest at the short end.}
\label{fig:gruniverse}
\end{figure}

\subsection{What would be wrong with copying?}

The oldest answer to a dark node is to copy a nearby smile --- take the
index smile at the same expiry, add a fixed spread, publish.  Copying
fails on three counts, and naming them tells us what the graph
must do instead.  First, it throws away the dark node's own
information: the sister name has a perfectly good smile from
yesterday, whose shape --- its skew, its wings, its own level --- is far
better than the index's shape with a spread.  Second, it has no notion
of \emph{how related} two nodes are: a copy is a copy, whether the
names move in lockstep or barely at all.  Third, it reports no
uncertainty: the copied smile arrives with the same face whether the
source is next door or three hops away.

The repairs, in order.  The dark node's own information is exactly
what \cref{ch:infer} already maintains: every node --- lit or dark ---
carries a \emph{baseline}, yesterday's posterior transported to
today's market by Chapter~9's rule, with a variance grown on
Chapter~8's clock.  This is the division of labour drawn at the end
of \cref{sec:infaudit}: persistence supplies each dark node's
baseline; what it cannot supply is this morning's \emph{news}.  So the
graph should not build dark smiles from scratch --- it should only move
them, and only as far as the lit evidence warrants.  Formally: the
unknown is not the dark node's handle but its \emph{innovation}, the
same departure-from-baseline that panel~(b) measures at the lit nodes.
Write $\Theta_i$ for that unknown at node $i$ --- capital theta, one
number per node for now.  If the morning's evidence says nothing about
node $i$, then $\Theta_i=0$: publish the baseline,
unchanged.  Everything the graph adds is a departure from zero that
the lit innovations must pay for.

The second repair --- how related --- is the subject of
\cref{sec:gredge,sec:grmeet}: relatedness becomes a stated amplitude
and a stated trust on each edge.  The third --- uncertainty --- runs
through the whole chapter: every filled node will report a variance,
the variance will widen with distance from the lit set, and
\cref{sec:grcomplete} runs Chapter~10's standardized-error test on
the published bands: the standardized errors $\mathcal{Z}$ of
\cref{eq:infz} must scatter with standard deviation one.

\subsection{The carrier, and the scope of one number}

One decision before the mathematics: \emph{what} moves along the
edges?  Not raw option prices (contracts differ node to node), and not
whole curves (a curve is not a number).  The reference implementation
propagates the smile's three ATM handles --- the level $\sigma(0)$, the
skew $s_0$, and the curvature (which, as in Chapter~10, keeps no
symbol and is named in words) --- the same compact, model-agnostic
state that \cref{sec:inffilter} filtered, and for the same reasons:
they are comparable across nodes, stable in meaning, and they retarget
into an arbitrage-free smile through Part~I's fits, so the
graph never publishes a raw curve.  Everything a smile is beyond its
three handles --- its wings above all --- simply rides the baseline:
the graph moves the belly and only the belly, the same division of
labour \cref{sec:infaudit} drew for persistence.  Each handle travels
the graph as its own scalar field with its own edge parameters.  The
chapter therefore develops everything for \emph{one} scalar field ---
$\Theta_i$, in volatility points --- and the staged morning is
constructed to move levels only (its skew and curvature fields carry
zero innovation everywhere and are not solved), so one field
suffices; where a handle-specific point matters, it is
flagged.  We now know
what the graph must infer.  The next question is what one edge is
allowed to say.
